% ── §1.3: Στόχος της Διπλωματικής ───────────────────────────

\section{Στόχος της Διπλωματικής}
\label{sec:objective}

Ο κεντρικός στόχος της παρούσας διπλωματικής εργασίας είναι ο σχεδιασμός
και η υλοποίηση ενός ολοκληρωμένου συστήματος Ομοσπονδιακής Μάθησης για
Αναγνώριση Δραστηριότητας σε Περιβάλλοντα Υποβοηθούμενης Διαβίωσης. Το
σύστημα αντιμετωπίζει ταυτόχρονα τα δύο θεμελιώδη ελλείμματα που
εντοπίστηκαν στη βιβλιογραφία (§~\ref{sec:motivation}): την έλλειψη
αξιόπιστης μέτρησης συνεισφοράς και την απουσία ολοκληρωμένου μηχανισμού
κινήτρων για ρεαλιστικά AAL σενάρια. Συγκεκριμένα, ενσωματώνει:

\begin{enumerate}
    \item \textbf{Μηχανισμό Κινήτρων:} Σχεδιασμός παιγνίου Stackelberg
    για την παρακίνηση ορθολογικών χρηστών να συνεισφέρουν ποιοτικά
    δεδομένα, με μαθηματικά αποδεδειγμένη σύγκλιση σε Nash Equilibrium.

    \item \textbf{Δίκαιη Αξιολόγηση Συνεισφοράς:} Υπολογισμός Shapley
    Values και μεθόδου Leave-One-Out (LOO) για τη μέτρηση της οριακής
    συνεισφοράς κάθε client στην αναγνωστική ικανότητα του καθολικού
    μοντέλου, με προσέγγιση μέσω gradient similarity για αποδοτική
    εκτέλεση.

    \item \textbf{Αρχιτεκτονική 1D-CNN:} Σχεδιασμός μονοδιάστατου
    Συνελικτικού Νευρωνικού Δικτύου ειδικά προσαρμοσμένου για την
    ανάλυση χρονοσειρών αισθητήρων HAR σε Non-IID συνθήκες FL.

    \item \textbf{Αξιολόγηση σε Ρεαλιστικές Συνθήκες:} Εκτίμηση
    απόδοσης τόσο υπό ελεγχόμενες (UCI HAR dataset) όσο και υπό
    πραγματικές συνθήκες (AAL dataset), με Non-IID κατανομή
    δεδομένων μεταξύ clients.
\end{enumerate}

Οι παραπάνω στόχοι δεν είναι ανεξάρτητοι: η αξιολόγηση συνεισφοράς τροφοδοτεί τον μηχανισμό κινήτρων, ο μηχανισμός κινήτρων επηρεάζει τη συμμετοχή και άρα την ποιότητα εκπαίδευσης του μοντέλου, και η απόδοση του μοντέλου επαληθεύει τη χρησιμότητα της αξιολόγησης — δημιουργώντας ένα κλειστό βρόχο σχεδιασμού που αξιολογείται ολιστικά στα πειράματα (Κεφάλαιο~6).
